%==============================================================================
% 6 个非凸源识别
%==============================================================================
\subsubsection{非凸性来源识别}

为使后续凸化推导可被严格审阅，本节预先识别原问题中的全部非凸源，并给出每项的数学依据。原始问题(P1)中的约束编号为(5a)--(5h)，详见§\ref{sec:problem}。

\begin{table}[htbp]
\centering
\caption{原问题(P1)的非凸源识别}\label{tab:nonconvex_sources}
\begin{tabular}{clll}
\toprule
编号 & 非凸源 & 受影响约束 & 凸性违背类型 \\
\midrule
NC1 & SINR 分式结构 & (5b), (5c) & 凸集对乘法不封闭 \\
NC2 & 半无限鲁棒约束 & (5b), (5c) & 无限约束违背凸性 \\
NC3 & 二进制组合约束 & (5h) & 离散集 $\{0,1\}$ 非凸 \\
NC4 & rank-1 隐含约束 & (5a)--(5f) & rank-1 集合非凸 \\
NC5 & 波束-信道双线性耦合 & (5b), (5d) & 双线性函数非凸 \\
NC6 & 矩阵逆约束 & (5d) & 逆映射保凸条件不满足 \\
\bottomrule
\end{tabular}
\end{table}

\textbf{非凸性的数学判据}：凸集 $\mathcal{C}$ 满足 $\forall \vect{x},\vect{y}\in\mathcal{C},\forall \theta\in[0,1]: \theta\vect{x}+(1-\theta)\vect{y}\in\mathcal{C}$。以此判据，下表给出关键约束的凸性状态：

\begin{itemize}
    \item $\|\vect{w}\|_2^2 \leq \alpha$：凸（二次锥）
    \item $\text{tr}(\mat{H}_k \mat{W}_k) \geq \beta$（$\mat{W}_k \succeq \mat{0}$）：凸（半正定锥 + 线性函数）
    \item $\text{rank}(\mat{W}_k) = 1$：\textbf{非凸}（rank-1 集不是仿射集）
    \item $\frac{x^2}{y} \leq \alpha, y > 0$：\textbf{非凸}（分式结构不保持凸性）
    \item $b \in \{0,1\}$：\textbf{非凸}（离散点集）
\end{itemize}

%==============================================================================
% 7 步凸化方法链
%==============================================================================
\subsubsection{七步凸化方法链}

针对上述 6 个非凸源，本节按编号给出 7 步凸化方法链。每步包括：受影响的非凸源、变换前后数学形式、变换性质（严格等价 / 紧致松弛 / 保守近似）、命题证明。

%------------------------------------------------------------------------------
\paragraph{Step 1：全局变量提升（消除 NC1, NC5）}
\textbf{变换性质}：\textbf{严格等价}。

将向量波束 $\vect{w}_k$ 提升为协方差矩阵：
\begin{equation}
\mat{W}_k = \vect{w}_k \vect{w}_k^H \in \mathbb{C}^{MN_t \times MN_t} \tag{15}
\end{equation}
\begin{equation}
\mat{Z} = \sum_p \vect{z}_p \vect{z}_p^H \in \mathbb{C}^{MN_t \times MN_t} \tag{16}
\end{equation}

\begin{proposition}[变量提升等价性]\label{prop:lifting}
当 $\mat{W}_k = \vect{w}_k \vect{w}_k^H$ 时，$\tr(\mat{H}_k \mat{W}_k) = |\vect{h}_k^H \vect{w}_k|^2$，且变量提升在 $\vect{w}_k \mapsto \mat{W}_k$ 上为一一映射。
\end{proposition}

\textbf{凸性影响}：目标函数 $\sum_k \|\vect{w}_k\|^2 = \sum_k \tr(\mat{W}_k)$ 由二次型转化为关于 $\mat{W}_k$ 的线性函数，\textbf{目标变凸}。

%------------------------------------------------------------------------------
\paragraph{Step 2：SDR 松弛（消除 NC4）}
\textbf{变换性质}：\textbf{紧致松弛}。$K\leq 2$ 时严格等价；$K>2$ 时为下界。

原 rank-1 约束 $\mat{W}_k = \vect{w}_k \vect{w}_k^H$ 等价于
\begin{equation}
\mat{W}_k \succeq \mat{0}, \quad \rank(\mat{W}_k) = 1
\end{equation}
SDR 丢弃 rank-1 约束，仅要求
\begin{equation}
\mat{W}_k \succeq \mat{0} \tag{S2.1}
\end{equation}
可行域由 $\mathcal{F}_{\text{rank-1}} = \{\mat{W}_k \succeq \mat{0} : \rank(\mat{W}_k)=1\}$ 扩大为 $\mathcal{F}_{\text{SDR}} = \{\mat{W}_k \succeq \mat{0}\}$（凸 PSD 锥）。

\begin{proposition}[SDR 紧致性条件]\label{prop:sdr_tight}
SDR 与原问题等价的充要条件是最优解满足 $\rank(\mat{W}_k^*) = 1$。
\begin{enumerate}
    \item $K \leq 2$ 时 SDR 紧致（由 Sidiropoulos, Davidson, Luo 2006 *IEEE Trans. SP* Theorem 1 关于 MISO 多播问题保证）。
    \item 高 SNR regime 下 SDR 近似紧致。
    \item 一般 $K>2$ 情形：高斯随机化（$\xi_l \sim \mathcal{CN}(\mat{0}, \mat{W}_k^*)$）以 $O(1/L)$ 性能损失提取可行 rank-1 解。
\end{enumerate}
\end{proposition}

\textbf{等价/上下界}：紧致时严格等价；一般情况为目标值下界（松弛扩大可行域 $\Rightarrow$ 最优值 $\leq$ 原值）。

%------------------------------------------------------------------------------
\paragraph{Step 3：S-Procedure 处理鲁棒 SINR（消除 NC2）}
\textbf{变换性质}：\textbf{保守近似}。用 1.75 dB 功率余量换 closed-form LMI 表示。

Worst-case SINR 约束为半无限优化：
\begin{equation}
\min_{\norm{\Delta\vect{h}_k} \leq \epsilon_h \norm{\hat{\vect{h}}_k}} \frac{|(\hat{\vect{h}}_k + \Delta\vect{h}_k)^H \mat{W}_k (\hat{\vect{h}}_k + \Delta\vect{h}_k)|}{\sum_{j\neq k}|(\hat{\vect{h}}_k + \Delta\vect{h}_k)^H \mat{W}_j(\hat{\vect{h}}_k + \Delta\vect{h}_k)| + \sigma_c^2} \geq \gamma_k
\end{equation}

\begin{lemma}[Cauchy-Schwarz 最坏情形]\label{lem:cauchy_wc}
对任意 $\vect{x}, \vect{y}$ 和 $\Delta\vect{y}$ 满足 $\|\Delta\vect{y}\| \leq \epsilon\|\vect{y}\|$：
\begin{equation}
\min_{\|\Delta\vect{y}\|\leq\epsilon\|\vect{y}\|} |\vect{x}^H(\vect{y}+\Delta\vect{y})|^2 = |\vect{x}^H\vect{y}|^2 (1-\epsilon)^2
\end{equation}
\textit{证明}：由 $| \vect{x}^H \Delta\vect{y}| \leq \|\vect{x}\| \cdot \|\Delta\vect{y}\| \leq \|\vect{x}\| \cdot \epsilon\|\vect{y}\|$，当 $\Delta\vect{y} = -\epsilon \frac{\|\vect{y}\|}{\|\vect{x}\|} \vect{x}$ 时取等。$\square$
\end{lemma}

应用引理~\ref{lem:cauchy_wc}（分子取下界、分母取上界）并交叉相乘，worst-case 鲁棒约束化为关于鲁棒门限 $\gamma_k^{\text{robust}}$ 的线性约束（详见式~\eqref{eq:lmi_comm}）。

\textbf{等价/上下界}：保守上界。真实最坏情形 SINR 高于此近似值 $\Rightarrow$ 求解出的可行域略大于真实鲁棒可行域，但保证真实信道下的鲁棒性。代价为约 1.75 dB 功率余量（$\epsilon_h=0.10$ 时 $\eta_h \approx 0.669$）。

%------------------------------------------------------------------------------
\paragraph{Step 4：SINR 分式线性化（消除 NC1 通信部分）}
\textbf{变换性质}：\textbf{严格等价}（前提：分母 $> 0$）。

在 Step 3 的鲁棒近似下，SINR 仍为分式形式。交叉相乘化为线性矩阵不等式：
\begin{equation}
\tr(\hat{\mat{H}}_k \mat{W}_k) - \gamma_k^{\text{robust}} \sum_{j\neq k} \tr(\hat{\mat{H}}_k \mat{W}_j) \geq \gamma_k^{\text{robust}} \sigma_c^2 \tag{S4.1}
\end{equation}
其中 $\hat{\mat{H}}_k = \hat{\vect{h}}_k \hat{\vect{h}}_k^H$。

\begin{proposition}[分式线性化等价性]\label{prop:frac}
当分母 $\sum_{j\neq k}|\hat{\vect{h}}_k^H \mat{W}_j \hat{\vect{h}}_k| + \sigma_c^2 > 0$ 时，
\begin{equation}
\frac{A}{B} \geq \gamma \iff A \geq \gamma B
\end{equation}
\end{proposition}

\textbf{凸性影响}：式~\eqref{eq:lmi_comm}（等价于式(S4.1)经 S-Procedure 提升后）成为关于 $\mat{W}_k$ 的 LMI，\textbf{通信约束变凸}。

%------------------------------------------------------------------------------
\paragraph{Step 5a：PCRB 仿射展开（消除 NC6）}
\textbf{变换性质}：\textbf{严格等价}（在 Assumption 1 下）。

Fisher 信息矩阵是 $\mat{R}_X$ 的仿射函数。记
\begin{equation}
\mat{F}_p = \frac{2}{\sigma_s^2} \Real\Big\{ \nabla_{\boldsymbol{\theta}_p}^H \vect{g}_p \cdot \nabla_{\boldsymbol{\theta}_p} \vect{g}_p^H \Big\} \in \mathbb{S}_+^{D\times D} \tag{13}
\end{equation}

PCRB 约束化为
\begin{equation}
\tr(\mat{F}_p \mat{R}_X) \geq \Gamma_{\text{Track},p} \tag{S5.1}
\end{equation}

\begin{proposition}[PCRB 线性化等价性]\label{prop:pcrb}
在 Assumption 1（$\nabla_{\boldsymbol{\theta}_p}\vect{g}_p$ 在当前时隙已知）下，$\tr(\mat{J}_p^{\text{data}}) = \tr(\mat{F}_p \mat{R}_X)$。
\end{proposition}

\textbf{凸性影响}：式(S5.1)关于 $\mat{R}_X$ 是线性约束，\textbf{感知 PCRB 约束变凸}。

%------------------------------------------------------------------------------
\paragraph{Step 5b：感知 SINR 线性化}
\textbf{变换性质}：\textbf{严格等价}（rank-1 最优解为匹配滤波）。

感知 SINR 约束 $\tr(\vect{g}_p\vect{g}_p^H \mat{Z}) \geq \gamma_S^{\text{PoD}} \sigma_s^2$ 中，$\vect{g}_p\vect{g}_p^H$ 为已知常数半正定矩阵，$\tr(\cdot)$ 关于 $\mat{Z}$ 为线性映射。

\begin{proposition}[感知 SINR 线性化等价性]\label{prop:sens_sinr}
当 $\mat{Z} = \vect{z}\vect{z}^H$（rank-1）时，最优感知波束为 $\vect{z}^* = \sqrt{P_S} \vect{g}_p/\|\vect{g}_p\|$（匹配滤波），且此时约束取等。
\end{proposition}

%------------------------------------------------------------------------------
\paragraph{Step 6：功率约束已为凸}
功率约束 $\sum_k \tr(\mat{W}_k) + \tr(\mat{Z}_m) \leq P_{\max}$ 关于 $\mat{W}_k, \mat{Z}_m$ 是线性约束，\textbf{已为凸}，无需变换。

%------------------------------------------------------------------------------
\paragraph{Step 7：AP 选择两步分解（处理 NC3）}
\textbf{变换性质}：\textbf{启发式下界}。外层离散 AP 组合为 NP-hard，内层连续变量在固定 AP 集合上为凸 SDP。

本工作采用两步分解：
\begin{enumerate}
    \item \textbf{外层（离散）}：基于大尺度衰落 $\text{PL}(d_{m,p})$ 排序选择 top-$N_{\text{req}}$ AP，确定 $\mathcal{M}^{\text{all}}$。
    \item \textbf{内层（凸）}：在固定 AP 集合上求解凸 SDP (P3)。
\end{enumerate}

\begin{remark}[AP 选择的工程取舍]\label{rem:ap_heuristic}
AP 选择问题本身是 NP-hard（$K$-medoid 问题的变种）。本工作以外层启发式 + 内层凸 SDP 的分解策略，以多项式复杂度获得工程可接受的次优解。复杂度下界证明详见配套文档~\cite{convex_chain} §6。
\end{remark}

%==============================================================================
% 凸化链总结表
%==============================================================================
\subsubsection{凸化链总结}

下表汇总 7 步凸化方法的变换类型与性能影响：

\begin{table}[htbp]
\centering
\caption{七步凸化方法链汇总}\label{tab:convex_chain}
\begin{tabular}{clll}
\toprule
Step & 变换 & 消除非凸源 & 变换类型 \\
\midrule
1 & 全局变量提升 $\vect{w}\vect{w}^H \to \mat{W}$ & NC1, NC5 & 严格等价 \\
2 & SDR 松弛（丢 rank-1） & NC4 & 紧致松弛（$K\leq 2$ 等价）\\
3 & S-Procedure 闭式界 & NC2 & 保守近似（1.75 dB 余量）\\
4 & SINR 分式线性化 & NC1（通信） & 严格等价 \\
5a & PCRB 仿射展开 & NC6 & 严格等价（Assumption 1）\\
5b & 感知 SINR 线性化 & NC1（感知） & 严格等价 \\
6 & 功率约束（已凸） & --- & 已是凸 \\
7 & AP 选择两步分解 & NC3 & 启发式下界 \\
\bottomrule
\end{tabular}
\end{table}

\textbf{关键结论}：通信-感知物理层的凸化（Step 1--6）几乎全部严格等价，仅 Step 3 为工程可接受的保守上界，Step 2 在 $K \leq 2$ 时紧致；AP 选择（Step 7）采用启发式以保证多项式复杂度。

%==============================================================================
% 最终凸问题 P3
%==============================================================================
\subsubsection{最终凸 SDP 问题 (P3)}

经 7 步凸化方法链，原问题(P1)化为以下\textbf{标准凸 SDP}：

\begin{align}
\min_{\{\mat{W}_k\}, \mat{Z}, \{\mu_k\}} \quad & \sum_{k=1}^K \tr(\mat{W}_k) + \tr(\mat{Z}) \tag{P3} \\
\text{s.t.} \quad & \begin{bmatrix} \mat{A}_k + \mu_k \mat{I} & \mat{A}_k \hat{\vect{h}}_k \\[1.5ex] \hat{\vect{h}}_k^H \mat{A}_k & \hat{\vect{h}}_k^H \mat{A}_k \hat{\vect{h}}_k - \sigma_c^2 - \mu_k \epsilon_h^2 \end{bmatrix} \succeq \mat{0}, \quad \forall k \label{eq:lmi_comm} \\
& \tr(\vect{g}_p \vect{g}_p^H \mat{Z}) \geq \gamma_S^{\text{PoD}} \sigma_s^2, \quad \forall p \label{eq:sinr_sens_sdp} \\
& \tr(\mat{F}_p \mat{R}_X) \geq \Gamma_{\text{Track},p}, \quad \forall p \label{eq:pcrb_sdp} \\
& \tr(\mat{E}_m \mat{R}_X) \leq P_{\max}, \quad \forall m \label{eq:power_sdp} \\
& \mat{W}_k \succeq \mat{0}, \quad \forall k \label{eq:psd_w} \\
& \mat{Z} \succeq \mat{0} \label{eq:psd_z} \\
& \mu_k \geq 0, \quad \forall k \label{eq:mu_nonneg}
\end{align}

其中 $\mat{A}_k = \frac{1}{\gamma_k} \mat{W}_k - \sum_{j \neq k} \mat{W}_j$，$\mat{F}_p$ 见式(13)，$\mat{E}_m$ 为 AP $m$ 的选择矩阵，$\mu_k \geq 0$ 为 S-Procedure 引入的松弛变量。

\begin{table}[htbp]
\centering
\caption{问题(P3)约束结构与变量对照}\label{tab:p3_structure}
\begin{tabular}{clll}
\toprule
约束编号 & 类型 & 涉及变量 & 物理意义 \\
\midrule
\eqref{eq:lmi_comm} & LMI ($MN_t+1$阶) & $\mat{W}_k, \mu_k$ & 通信最坏情况鲁棒性 \\
\eqref{eq:sinr_sens_sdp} & 线性不等式 & $\mat{Z}$ & 感知探测概率 \\
\eqref{eq:pcrb_sdp} & 线性不等式 & $\mat{R}_X$ & 跟踪精度(PCRB) \\
\eqref{eq:power_sdp} & 线性不等式 & $\mat{R}_X$ & 单AP峰值功率 \\
\eqref{eq:psd_w}--\eqref{eq:psd_z} & 半正定锥 & $\mat{W}_k, \mat{Z}$ & SDR松弛 \\
\eqref{eq:mu_nonneg} & 非负象限 & $\mu_k$ & S-Procedure松弛 \\
\bottomrule
\end{tabular}
\end{table}

\begin{theorem}[问题凸性]\label{thm:convex_p3}
问题(P3)是标准凸 SDP，目标函数为线性、约束均为线性矩阵不等式（LMI）或半正定锥约束，可由 MOSEK / SeDuMi / SDPT3 等标准 SDP 求解器在多项式时间内求解到任意精度。求解复杂度为 $O((MN_t)^{3.5})$，对 $M=16, N_t=4$ 约 5--10 秒。
\end{theorem}

\begin{theorem}[SDR 紧致性]\label{thm:sdr_tight}
对于总功率最小化问题，当 $K \leq 2$ 时，SDR 紧致，即最优解满足 $\rank(\mat{W}_k^*) = 1$。高 SNR regime 下 SDR 近似紧致。一般 $K>2$ 情形，通过高斯随机化恢复波束，性能损失上界为 $O(1/L)$（$L$ 为候选数）。
\end{theorem}

%==============================================================================
% 复杂度对比
%==============================================================================
\subsubsection{凸化前后的复杂度对比}

\begin{table}[htbp]
\centering
\caption{凸化前后问题形式与求解复杂度对比}\label{tab:complexity_compare}
\begin{tabular}{llll}
\toprule
问题 & 形式 & 求解器 & 复杂度 \\
\midrule
(P1) 原问题 & MINLP（混合整数非凸） & 无通用算法 & NP-hard \\
(P3) 凸 SDP & 凸半定规划 & MOSEK / SeDuMi & $O((MN_t)^{3.5})$ \\
\bottomrule
\end{tabular}
\end{table}

凸化将 NP-hard 的混合整数非凸问题转化为多项式时间可解的标准凸 SDP，这是凸化方法的核心价值。
